\section{The Grammar of Seeing and the Vocabulary of the Grave}
\label{sec:grammar}

The single Greek participle around which most modern speculation turns---rendered ``wrapped up'' in the Douay--Rheims, ``wrapped together'' in the King James Version, and ``folded'' in several late-twentieth-century translations---appears only three times in the entire New Testament. Its other two occurrences describe the burial of Jesus on Good Friday afternoon: St. Matthew writes that Joseph of Arimathea took the body and ``wrapped it up'' (\latin{enetylixen}) in a clean linen cloth (Matt 27:59), and St. Luke uses the exact same verb when Joseph takes the body down from the cross and ``wrapped him'' (\latin{enetylixen}) in fine linen (Luke 23:53). Before a single commentary is opened, the Greek concordance itself removes the domestic crease-folding of a dinner napkin from the vocabulary of John 20:7.

To see what St. John actually sets before his readers, every concrete noun, verb, and spatial phrase in John 20:3--10 must be read alongside its New Testament controls and the Clementine Vulgate. The text is extraordinarily spare. It contains no adjective describing the fabric's quality, no explanation of how the cloths came to rest where they did, and no narrator's aside about domestic customs. Every ounce of weight rests on the exact Greek words chosen.\endnote{Greek wording and critical variant notes throughout are checked directly against the Robinson--Pierpont Byzantine Textform (RP2018 Unicode CSV, which marks Nestle--Aland divergences inline), the SBL Greek New Testament (Michael W. Holmes, v1.2, for Luke), and MorphGNT SBLGNT v6.12 for complete New Testament lemma counts. Lexical definitions are verified against Joseph Henry Thayer, \emph{A Greek-English Lexicon of the New Testament}, corrected edition (New York: American Book Company, undated title page; 1889).}

\subsection{The Linen Cloths: \emph{Othonia}, \emph{Sind\={o}n}, and \emph{Keiriai}}

The Gospels use three distinct Greek nouns for the textiles of Jewish burial, and the difference among them explains how John 19--20 fits beside both the Synoptic burial narratives and the raising of Lazarus in John 11.

In John 20:5, 6, and 7, the disciples see \latin{ta othonia} (the plural of \latin{othonion}, rendered in the Vulgate as \latin{linteamina}, ``linen cloths''). Across the entire Greek New Testament, \latin{othonion} occurs only five times, and every single occurrence belongs to the burial and empty tomb of Jesus: once at the preparation of the body in John 19:40, three times inside the tomb in John 20:5--7, and once in Peter's visit in Luke 24:12. Thayer defines \latin{othonion} as a diminutive of \latin{othon\={e}} (a large linen sheet, which appears twice in Acts 10:11 and 11:5 for the great sheet let down from heaven in Peter's vision at Joppa): ``a piece of linen, small linen cloth: plur. strips of linen cloth for swathing the dead,'' noting that in Greek writers the word denotes linen bandages for wounds, ships' sails, and similar linen fabrics, and in the Septuagint renders Hebrew \latin{sadin} (Judg 14:13).

St. John prepares the reader for these plural cloths in the preceding chapter:
\begin{quote}
``And Nicodemus also came (he who at the first came to Jesus by night), bringing a mixture of myrrh and aloes, about an hundred pound weight. They took therefore the body of Jesus and bound it in linen cloths, with the spices, as the manner of the Jews is to bury.'' (John 19:39--40)
\end{quote}
The verb in John 19:40 is \latin{ed\={e}san} (``they bound,'' Vulgate \latin{ligaverunt}), governing \latin{othoniois} (``with linen cloths,'' Vulgate \latin{linteis}, with the Byzantine text adding the preposition \latin{en}, ``in''). The hundred Roman pounds of myrrh and aloes---roughly seventy-two modern pounds of fragrant resinous powder and gum---were packed within and between the linen wrappings as Joseph and Nicodemus bound the body hastily before the Sabbath began (John 19:42).

The Synoptic Gospels, by contrast, speak in the singular of a \latin{sind\={o}n} (fine linen sheet): Joseph wraps the body in ``a clean linen cloth'' (\latin{sindoni kathara}, Matt 27:59), buys ``fine linen'' and wraps him in it (\latin{eneil\={e}sen t\={e} sindoni}, Mark 15:46), and ``wrapped him in fine linen'' (\latin{enetylixen auto sindoni}, Luke 23:53). Elsewhere in the New Testament, \latin{sind\={o}n} appears only in Mark 14:51--52 for the single linen garment worn by the young man who fled naked in Gethsemane---six occurrences in total.

Already in the early fifth century, St. Augustine addressed the apparent tension between the Synoptics' single \latin{sind\={o}n} and John's plural \latin{othonia}. In \emph{The Harmony of the Gospels} (III.23), Augustine observes first that the Synoptics' mention of Joseph does not exclude Nicodemus from bringing additional linen fabrics, and second that even if only one principal shroud were involved, the accompanying head-cloth and binding strips were all flaxen textiles: ``when we take into account the handkerchief which was used for the head, and the bandages with which the whole body was swathed, and consider that all these were made of linen, we can see how, even although there was really but a single linen cloth there, it could still have been stated with the most perfect truth that `they wound Him in linen clothes.' For the phrase, linen clothes, is one applied generally to all textures made of flax.'' St. Luke himself confirms that no contradiction exists between the two terms: in Luke 23:53 Joseph wraps Jesus in a \latin{sind\={o}n}, and twelve verses later in Luke 24:12 Peter stoops down and sees \latin{ta othonia}, the linen cloths.

A third burial noun appears only once in the New Testament, at the tomb of Lazarus: \latin{keiriai} (``winding bands,'' Vulgate \latin{institis}), binding Lazarus's feet and hands (John 11:44). Unlike the broader term \latin{othonia}, \latin{keiriai} designates narrow binding straps or cords wound tightly around limbs so that the dead man emerges immobilized until bystanders loose him. Yet one noun bridges both tombs identically: the \latin{soudarion}.

\subsection{The Head-Cloth: \emph{Soudarion} and the History of English ``Napkin''}

In John 20:7, Simon Peter sees not only the linen cloths lying, ``and the napkin that had been about his head'' (\latin{kai to soudarion ho \={e}n epi t\={e}s kephal\={e}s autou}; Vulgate \latin{et sudarium, quod fuerat super caput ejus}).

The Greek word \latin{soudarion} is not a native Greek culinary term and not a Hebrew domestic word. It is a direct transliteration of the Latin \latin{sudarium}, derived from \latin{sudor} (``sweat''). Thayer defines \latin{soudarion} as ``a handkerchief, i.e. a cloth for wiping the perspiration from the face and for cleaning the nose: Lk. xix. 20; Acts xix. 12; also used in swathing the head of a corpse, Jn. xi. 44; xx. 7.'' The noun occurs exactly four times in the entire Greek New Testament, and not once in a dining room:

\begin{itemize}
\item \textbf{Luke 19:20.} In the parable of the ten pounds, the fearful servant says: ``Lord, behold here is thy pound, which I have kept laid up in a napkin'' (\latin{en soudari\={o}}, Vulgate \latin{in sudario}). He tied his master's silver coin inside a personal kerchief or sweat-cloth for safekeeping.
\item \textbf{Acts 19:12.} At Ephesus, miracles are worked through St. Paul so that ``there were brought from his body to the sick, handkerchiefs and aprons'' (\latin{soudaria \={e} simikinthia}, Vulgate \latin{sudaria et semicinctia}). Here \latin{soudaria} stands beside \latin{simikinthia} (another Latin loanword, \latin{semicinctia}, workmen's aprons)---sweat-cloths worn by a tentmaker at his trade.
\item \textbf{John 11:44.} When Jesus calls Lazarus from the grave, ``he that had been dead came forth, bound feet and hands with winding bands. And his face was bound about with a napkin'' (\latin{kai h\={e} opsis autou soudari\={o} periededeto}, Vulgate \latin{et facies illius sudario erat ligata}). Here the verb \latin{peride\={o}} (a New Testament \latin{hapax legomenon}, ``to bind round about'') describes a cloth tied around a dead man's face or jaw.
\item \textbf{John 20:7.} Inside the sepulchre of Jesus, Peter sees ``the napkin that had been about his head'' (\latin{epi t\={e}s kephal\={e}s autou}, literally ``upon his head''), lying apart from the \latin{othonia}.
\end{itemize}

Why, then, do English readers meet the word ``napkin'' in John 20:7? The answer is a straightforward chapter in the history of English. When the Catholic translators at Rheims rendered the Latin \latin{sudarium} in 1582 (followed in Challoner's Douay--Rheims revision) and when the King James translators rendered the Greek \latin{soudarion} in 1611, Early Modern English ``napkin'' (from Old French \latin{nappe}, a cloth, plus the diminutive \latin{-kin}) meant a small square of linen carried on the person---a handkerchief, sweat-cloth, or head-cloth---just as readily as a table cloth.\endnote{Thus both the Douay--Rheims and the King James Version use ``napkin'' identically in Luke 19:20 (storing a coin), John 11:44 (binding Lazarus's dead face), and John 20:7 (the cloth on Jesus' head), while shifting to ``handkerchiefs'' only in Acts 19:12. The Revised Version of 1895 retained ``napkin'' in Luke 19:20, John 11:44, and John 20:7, whereas modern public-domain versions such as the World English Bible (Catholic Edition) render ``cloth'' in John 11:44 and 20:7 and ``handkerchief'' in Luke 19:20 and Acts 19:12.} In the centuries since 1611, ordinary English narrowed ``napkin'' almost exclusively to the dining table, handing the personal sweat-cloth over to ``handkerchief.'' A modern reader who encounters ``the napkin that had been about his head'' without checking Luke 19:20, John 11:44, or Acts 19:12 easily pictures a dinner napkin where St. John wrote of a burial face-cloth.

\subsection{Rolled Up Apart: \emph{Entetyligmenon}, \emph{Ch\={o}ris}, and \emph{Eis Hena Topon}}

Verse 7 specifies the exact condition and position of the \latin{soudarion} in two contrasting clauses:
\begin{quote}
\latin{ou meta t\={o}n othoni\={o}n keimenon, alla ch\={o}ris entetyligmenon eis hena topon.}\\
\latin{non cum linteaminibus positum, sed separatim involutum in unum locum.}
\end{quote}
Four details govern the sense of this clause.

First, the participle \latin{entetyligmenon} is the perfect passive participle of \latin{entyliss\={o}}. Thayer gives its two senses with exact citations: ``to roll in, wrap in'' (as in Matt 27:59 and Luke 23:53, \latin{tina sindoni}, wrapping a body in a linen sheet) and ``to roll up, wrap together'' (passive in John 20:7, citing Aristophanes and Athenaeus). The Greek verb is built on \latin{tyliss\={o}} (to roll, coil, twine around). It does not mean creasing a square of fabric flat into crisp rectangular folds. When St. Luke wishes to describe rolling or folding up a synagogue book-roll in Luke 4:20, he uses \latin{ptyss\={o}} (a distinct verb occurring once in the New Testament, paired with \latin{anaptyss\={o}} in Luke 4:17). St. Jerome caught the Greek sense accurately in the Vulgate: \latin{involutum} (``rolled up,'' ``wrapped round''), the same Latin root (\latin{involvit}) he used in Matthew 27:59, Mark 15:46, and Luke 23:53 for Joseph wrapping the Lord's body in the shroud. Likewise, Douay--Rheims translates ``wrapped up,'' the King James Version ``wrapped together,'' and the Revised Version of 1895 ``rolled up.''

Second, the tense and voice of \latin{entetyligmenon} are perfect passive. In Greek syntax, the perfect participle depicts an existing state resulting from a completed action: when Peter enters, the head-cloth is already in a rolled-up or coiled state. The passive voice names the state of the cloth without specifying the agent or the mechanics of how it came into that state. That grammatical restraint leaves room both for St. John Chrysostom, who takes \latin{entetyligmenon} as the unhurried work of a hand setting the cloth aside, and for Henry Latham in 1901, who takes the same perfect passive participle as the cloth retaining the coiled shape it had been wrapped into around the Lord's head on Friday evening.

Third, the adverb \latin{ch\={o}ris} (``apart,'' ``separately,'' Vulgate \latin{separatim}) stands in direct antithesis to \latin{ou meta t\={o}n othoni\={o}n keimenon} (``not lying with the linen cloths''). In the English Revised Version of 1895 (``rolled up in a place by itself'') and the King James Version (``wrapped together in a place by itself''), \latin{ch\={o}ris} is folded into ``by itself''; the Douay--Rheims (``apart, wrapped up into one place'') preserves both \latin{ch\={o}ris} and \latin{eis hena topon} explicitly, faithfully mirroring the Vulgate's \latin{separatim involutum in unum locum}.

Fourth, the prepositional phrase \latin{eis hena topon} (literally ``into one place,'' idiomatically ``in one place by itself'') conveys collectedness and order. In late Greek, \latin{eis} frequently overlaps with \latin{en} (``in'') while still carrying a faint note of placement---having been gathered or set into a single spot of its own, rather than strewn across the tomb floor.

\subsection{Who Saw What: Three Verbs of Seeing and Two Disciples}

English translations smooth out a deliberate progression in St. John's narration of how the two Apostles approach the tomb. In Douay--Rheims, the same English verb ``saw'' translates three different Greek verbs across verses 5, 6, and 8:

\begin{enumerate}
\item \textbf{Verse 5 (\latin{blepei}):} The other disciple outruns Peter, reaches the entrance first, stoops down (\latin{parakypsas}, the same participle Luke 24:12 uses of Peter), and ``sees'' (\latin{blepei}, historic present) the linen cloths lying (\latin{keimena ta othonia}), yet does not go in. The verb \latin{blep\={o}} denotes ordinary visual perception---catching sight of what is visible from the low doorway.
\item \textbf{Verse 6 (\latin{the\={o}rei}):} Simon Peter arrives following him, goes straight into the sepulchre, and ``beholds'' or ``observes'' (\latin{the\={o}rei}, historic present) the linen cloths lying (\latin{ta othonia keimena}), and the \latin{soudarion} rolled up apart in one place (v. 7). The verb \latin{the\={o}re\={o}} frequently denotes sustained, attentive looking---taking in details one after another. Crucially, the entire description of the \latin{soudarion} in verse 7 hangs grammatically on Peter's \latin{the\={o}rei}. The younger disciple saw only the \latin{othonia} from outside; Peter enters first and sees both the \latin{othonia} and the separate, rolled-up \latin{soudarion}.
\item \textbf{Verse 8 (\latin{eiden}):} ``Then that other disciple also went in, who came first to the sepulchre: and he saw and believed'' (\latin{kai eiden kai episteusen}). Both verbs shift abruptly from the vivid historic present (\latin{blepei}, \latin{the\={o}rei}) into the aorist (\latin{eiden}, \latin{episteusen}).
\end{enumerate}

Brooke Foss Westcott weighed both the lexical choice and the tense shift carefully in his Greek commentary on St. John (1908, vol. 2, p. 339). Of verse 5 he notes that ``the simple sight here is distinguished from the intent regard of St. Peter when he entered the sepulchre; and in this connexion it is significant that St. John does not see the napkin, the small cloth, lying apart.'' Of verse 6 he adds: ``The abrupt change of tense marks a break in the progress of the thought. The entrance is courageously made: then follows the experience. The word expresses the earnest intent gaze of the apostle as his eye passes from point to point.''

At the same time, good Greek exegesis refuses to press lexical distinctions farther than St. John's own style permits. Across the Fourth Gospel, St. John often varies synonyms (\latin{blep\={o}}, \latin{the\={o}re\={o}}, \latin{hora\={o}}) for stylistic fullness without confining each verb to an airtight technical sense; in verse 12, Mary Magdalen uses the very same \latin{the\={o}rei} when she sees the two angels, and in verse 14 when she sees Jesus standing without recognizing him. The undeniable progression in John 20:5--8 does not depend on over-subtilizing Greek synonyms alone: it is built right into the physical sequence of actions. From outside the doorway, the younger disciple sees only the linen cloths; stepping inside, Peter inspects both the linen cloths and the separate head-cloth; stepping inside after Peter, the younger disciple sees what Peter has seen, and the Evangelist writes---in the singular---\latin{kai eiden kai episteusen}, ``and he saw and believed.''

Then arrives the sentence that divided the Fathers: \latin{oudep\={o} gar \={e}deisan t\={e}n graph\={e}n, hoti dei auton ek nekr\={o}n anast\={e}nai} (``For as yet they knew not the scripture, that he must rise again from the dead,'' v. 9). The grammatical number shifts abruptly between the two verses. Verse 8 is singular (\latin{eiden}, \latin{episteusen}: the beloved disciple); verse 9 is plural (\latin{\={e}deisan}: ``they knew not''), encompassing both Peter and the beloved disciple, and verse 10 remains plural (\latin{ap\={e}lthon}: ``the disciples therefore departed again to their home'').\endnote{In John 20:10, the Byzantine textform reads \latin{pros heautous} (reflexive, ``to themselves'' / ``to their own home,'' Vulgate \latin{ad semetipsos}), whereas the Nestle--Aland text reads \latin{pros autous} (or \latin{pros hautous}, carrying the same idiomatic sense ``to their lodging / home''). In John 20:6, the Nestle--Aland text includes \latin{kai} before \latin{Sim\={o}n Petros} (``Simon Peter also cometh''), which the Byzantine text omits.}

\Needspace{14\baselineskip}
\begin{fourstable}{Greek Expression (John 20:5--7)}{Clementine Vulgate}{Douay--Rheims / KJV (1582--1611)}{RV (1895) / WEB-C}
ta othonia keimena & linteamina posita & DR: the linen cloths lying; KJV: the linen clothes lying & RV: the linen cloths lying; WEB-C: the linen cloths lying \\
\midrule
to soudarion & sudarium & DR \& KJV: the napkin & RV: the napkin; WEB-C: the cloth \\
\midrule
entetyligmenon & involutum & DR: wrapped up; KJV: wrapped together & RV \& WEB-C: rolled up \\
\midrule
ch\={o}ris \dots\ eis hena topon & separatim \dots\ in unum locum & DR: apart \dots\ into one place; KJV: in a place by itself & RV \& WEB-C: in a place by itself \\
\end{fourstable}

Even before turning to St. Luke or the Fathers, this comparison (expanded across every New Testament parallel in Appendix~\ref{app:register}) settles the floor of the inquiry. John 20:6--7 describes burial cloths (\latin{othonia}) bound with a hundred pounds of spices and a burial face-cloth (\latin{soudarion}) rolled or wrapped up (\latin{entetyligmenon}) in a place apart (\latin{ch\={o}ris eis hena topon})---observed first by Simon Peter upon entering the tomb, and then by the beloved disciple.
